\topic{16}{Read errors, equation residuals and recorded times}
The calculated centerlines are compared at the actual Ghia reference sample
locations. Let $s_j$ be sample position $j$, $f_{\rm ref,j}$ its tabulated
velocity, and $f_{\rm num}(s_j)$ the computed profile interpolated linearly
to that position. The symbol $f$ stands for either the horizontal or vertical
velocity component in its corresponding centerline profile.

Define the relative profile error and largest absolute profile error by
\[
E_{{\rm rel},2}=\frac{\sqrt{\sum_j[f_{\rm num}(s_j)-f_{\rm ref,j}]^2}}
{\sqrt{\sum_j f_{\rm ref,j}^{\,2}}},\qquad
E_\infty=\max_j|f_{\rm num}(s_j)-f_{\rm ref,j}|.
\]
The sums and maximum run over the stated reference samples. A relative error
of $0.02$ is $2\%$, since percentage error is $100E_{{\rm rel},2}$.
The denominator above is nonzero for these Ghia profiles. The qualified
$Re=400$ vertical-velocity score excludes the separately marked published
sample at $x/L=0.9063$; the all-point score is also retained for comparison.

\begin{center}\begin{tabular}{lp{.68\linewidth}}\toprule
Reported quantity & What it measures\\\midrule
Profile error & Difference from independent velocity samples.\\
Momentum defect & Error in satisfying the discrete momentum equations.\\
Cell divergence & Error in satisfying the discrete cell continuity equations.\\
Iterations / steps & Completed steady updates or physical marching steps; their work differs.\\
Solver-call seconds & Recorded elapsed time for one solver call on the stated host.\\\bottomrule
\end{tabular}\end{center}
A small equation residual does not guarantee a small profile error. Mesh and
spatial approximations remain relevant after a solver has converged. The
streamfunction--vorticity method has its own update and steady-vorticity
checks; these are not finite-volume momentum residuals.

Axis labels such as $u/U$ and $p/(\rho U^2)$ denote the dimensional ratios
used to define the nondimensional fields in Section 1. The plotted stored
values already represent those ratios; they are not divided by the scales
a second time.

For timing comparisons, distinguish retained records from newly executed
cases. Match spacing, equations and stopping requirements before interpreting
seconds. Timing one run on one host does not establish a universal fastest
method.
